\section{Content-factor extraction and fixed-determinant preservation}
\label{sec:determinant-law}

The resolved formula \eqref{eq:fully-resolved-inner-sum} still uses
the value \(U_{\theta,b}^\ast(z)\).  The defining congruence forces a
factor \(B_{\theta,b}\) from that value.  Extracting it is essential
before invoking the fixed-determinant local theory of TPC-90.

\subsection{The two-stage determinant law}

On a soluble branch, put \(B=B_{\theta,b}\),
\(\tau=\tau_{\theta,b}\), and define
\begin{equation}
 V_{\theta,b}(z)
 =
 \frac{U_\theta(\tau+Bz)}{B}.
\label{eq:det-reduced-V}
\end{equation}
This is integral because \(B\mid U_\theta(\tau)\).  More explicitly,
\begin{align}
 \mathcal D_{\theta,b}(z)
 &:=
 D_\theta(\tau+Bz)
 =
 D_\theta(\tau)+\sigma_\theta Bz,
\label{eq:det-reduced-D}\\
 V_{\theta,b}(z)
 &=
 \frac{U_\theta(\tau)}B+a_\theta z.
\label{eq:det-reduced-V-affine}
\end{align}

\begin{theorem}[Content-progression determinant law]
\label{thm:content-progression-determinant}
Let a nonempty literal resolved term satisfy
\[
 b=c\kappa\mid P_\theta,\qquad
 b\mid\sigma_\theta U_\theta(t),
\qquad
 B=\frac b{(b,\sigma_\theta)}.
\]
Then:
\begin{enumerate}[label=(\roman*),leftmargin=2.2em]
\item the unreduced reparametrized pair
\[
 \mathcal D_{\theta,b}(z),\qquad
 U_\theta(\tau+Bz)
\]
has coefficient determinant \(Bh_0\);
\item the second value factors exactly as
\[
 U_\theta(\tau+Bz)=B\,V_{\theta,b}(z);
\]
\item the reduced pair
\[
 \mathcal D_{\theta,b}(z)
 =
 D_\theta(\tau)+(\sigma_\theta B)z,\qquad
 V_{\theta,b}(z)
 =
 \frac{U_\theta(\tau)}B+a_\theta z
\]
has coefficient determinant exactly \(h_0\):
\begin{equation}
 \boxed{
 (\sigma_\theta B)\frac{U_\theta(\tau)}B
 -
 a_\theta D_\theta(\tau)
 =
 h_0;}
\label{eq:det-preserved}
\end{equation}
\item on every nonempty literal branch,
\begin{equation}
 (a_\theta,\sigma_\theta B)=1,\qquad
 (a_\theta\sigma_\theta B,h_0)=1;
\label{eq:det-reduced-slope-primitivity}
\end{equation}
\item the reduced forms are individually primitive:
\begin{equation}
 \gcd(D_\theta(\tau),\sigma_\theta B)=1,\qquad
 \gcd(U_\theta(\tau)/B,a_\theta)=1.
\label{eq:det-reduced-form-primitivity}
\end{equation}
\end{enumerate}
\end{theorem}

\begin{proof}
The unreduced forms have origins
\(D_\theta(\tau),U_\theta(\tau)\) and slopes
\(\sigma_\theta B,a_\theta B\).  Their determinant is
\[
 (\sigma_\theta B)U_\theta(\tau)
 -(a_\theta B)D_\theta(\tau)
 =
 B\{\sigma_\theta U_\theta(\tau)
 -a_\theta D_\theta(\tau)\}
 =
 Bh_0.
\]
The progression was chosen so that \(B\mid U_\theta(\tau)\), proving
the factorization.  Dividing the second origin and slope by \(B\)
gives \eqref{eq:det-preserved}.

For a nonempty term, \(b\) divides both physical targets.  Primitive
target support therefore gives
\begin{equation}
 (b,h_0)=1
\label{eq:det-b-coprime-h}
\end{equation}
and, on every contributing source atom, \(b\) is coprime to the
moving row and \(j_\theta\).  In particular \((B,h_0)=1\).
Solubility gives \((a_\theta,B)=1\) by
\cref{lem:one-progression}; combine this with
\((a_\theta,\sigma_\theta)=1\) and
\((a_\theta\sigma_\theta,h_0)=1\) from
\eqref{eq:affine-native-determinant} to obtain
\eqref{eq:det-reduced-slope-primitivity}.

Any common divisor of \(D_\theta(\tau)\) and
\(\sigma_\theta B\) divides the left side of
\eqref{eq:det-preserved}, hence \(h_0\); this contradicts
\((\sigma_\theta B,h_0)=1\).  The second primitive identity is
proved in the same way using \((a_\theta,h_0)=1\).
\end{proof}

\subsection{Exact M\"obius extraction}

\begin{lemma}[Forced content factor in the M\"obius sign]
\label{lem:forced-mobius-factor}
For \(B\ge1\) and \(V\ge1\),
\begin{equation}
 \mu(BV)
 =
 \begin{cases}
 \mu(B)\mu(V)\one_{(B,V)=1},&\mu^2(B)=1,\\
 0,&\mu^2(B)=0.
 \end{cases}
\label{eq:forced-mobius-factor}
\end{equation}
\end{lemma}

\begin{proof}
If \(B\) is not squarefree, then \(BV\) is not squarefree.  If \(B\)
is squarefree and \((B,V)>1\), again \(BV\) has a squared prime
factor.  On the remaining branch, multiplicativity of \(\mu\) gives
the identity.
\end{proof}

\begin{theorem}[Fixed-\(h_0\) reduced native export]
\label{thm:fixed-h-reduced-export}
In the fully resolved low-window formula, every branch with
\(\mu^2(B_{\theta,b})=0\) vanishes.  On the remaining branches,
\eqref{eq:fully-resolved-inner-sum} equals
\begin{align}
 \mu(B_{\theta,b})\zeta_{\theta,b,c,r}
 \sum_{z\in I_{\theta,b}^\ast}
 &\mu(\mathcal D_{\theta,b}(z))
 \mu(V_{\theta,b}(z))
 \one_{(B_{\theta,b},V_{\theta,b}(z))=1}
\notag\\
 &\times
 \chi_\theta(\tau_{\theta,b}+B_{\theta,b}z)
 W_{\theta,X}(\tau_{\theta,b}+B_{\theta,b}z)
 \notag\\
 &\times
 e_{cq_X}(-r\varepsilon_\theta\Omega_{\theta,b}z).
\label{eq:fixed-h-reduced-inner}
\end{align}
The two M\"obius arguments in
\eqref{eq:fixed-h-reduced-inner} are individually primitive affine
forms with the same prescribed determinant \(h_0\).  All extracted
factors and masks remain literal.
\end{theorem}

\begin{proof}
Use
\[
 U_{\theta,b}^\ast(z)
 =
 B_{\theta,b}V_{\theta,b}(z)
\]
in \eqref{eq:fully-resolved-inner-sum} and apply
\cref{lem:forced-mobius-factor}.  The determinant and primitive
claims are \cref{thm:content-progression-determinant}.
\end{proof}

\begin{corollary}[Exact TPC-90 local split after content extraction]
\label{cor:tpc90-after-extraction}
For every reduced branch in
\cref{thm:fixed-h-reduced-export}, the TPC-90 identity applies:
\begin{align}
 &\mu(\mathcal D_{\theta,b}(z))
 \mu(V_{\theta,b}(z))
\notag\\
 &\quad=
 \mu^2(\mathcal D_{\theta,b}(z))
 \mu^2(V_{\theta,b}(z))
 \lambda(\mathcal D_{\theta,b}(z))
 \lambda(V_{\theta,b}(z)).
\label{eq:tpc90-reduced-split}
\end{align}
Its simultaneous squarefree local root counts are those for a
primitive determinant-\(h_0\) pair.  The extra mask
\(\one_{(B_{\theta,b},V_{\theta,b}(z))=1}\), the sign
\(\mu(B_{\theta,b})\), and the additive determinant character remain
outside that identity.
\end{corollary}

\begin{proof}
Apply the exact identity \(\mu=\mu^2\lambda\) to the two positive
reduced forms.  The hypotheses of the TPC-90 primitive local
classification follow from
\cref{thm:content-progression-determinant}.
\end{proof}

\subsection{The precise provenance boundary}

We now isolate the exact obstruction to an untracked determinant
change.

\begin{theorem}[No silent determinant relabeling]
\label{thm:no-silent-relabeling}
After the substitution \(t=\tau+Bz\), the unreduced pair has
determinant \(Bh_0\).  Therefore it cannot be entered into a theorem
whose hypothesis is determinant \(h_0\) by simply retaining its two
unreduced values and changing the determinant label.  A valid
fixed-\(h_0\) entry must instead perform
\[
 U(\tau+Bz)=B\,V(z)
\]
and retain:
\begin{enumerate}[label=(\alph*),leftmargin=2.2em]
\item the branch deletion \(\mu^2(B)=0\);
\item the outer sign \(\mu(B)\);
\item the coprimality mask \(\one_{(B,V(z))=1}\); and
\item the original content, phase, polarization, and outer
provenance keys.
\end{enumerate}
Dropping any one of (a)--(c) changes the literal coefficient.
\end{theorem}

\begin{proof}
The determinant statements follow from
\cref{thm:content-progression-determinant}.  The three coefficient
effects are the exhaustive cases of
\eqref{eq:forced-mobius-factor}.  The remaining provenance is needed
to reverse the finite expansion in
\cref{thm:content-expanded-window}; without it one no longer has the
source coefficient \eqref{eq:row-literal-A}.
\end{proof}

\begin{example}[The two determinants are visibly different]
\label{ex:two-determinants}
Take
\[
 D(t)=t,\qquad U(t)=t+1,\qquad h_0=1.
\]
Restrict to \(2\mid U(t)\), so \(t=1+2z\).  The unreduced pair
\[
 1+2z,\qquad2+2z
\]
has determinant \(2\).  Extracting the forced factor from the second
value gives
\[
 1+2z,\qquad1+z,
\]
whose determinant is \(1\).  The equality
\[
 \mu(2+2z)
 =
 -\mu(1+z)\one_{(2,1+z)=1}
\]
shows why the sign and coprimality mask cannot be omitted.  This is a
finite algebraic example, not a model of the complete physical
packet.
\end{example}
